% !TEX root =
% ============================================================================
%  Background chapter for the MSc thesis.
%  Self-contained: requires amsmath, amssymb, and a natbib-style \citep/\citet.
%  Drop into the thesis with % !TEX root =
% ============================================================================
%  Background chapter for the MSc thesis.
%  Self-contained: requires amsmath, amssymb, and a natbib-style \citep/\citet.
%  Drop into the thesis with % !TEX root =
% ============================================================================
%  Background chapter for the MSc thesis.
%  Self-contained: requires amsmath, amssymb, and a natbib-style \citep/\citet.
%  Drop into the thesis with % !TEX root =
% ============================================================================
%  Background chapter for the MSc thesis.
%  Self-contained: requires amsmath, amssymb, and a natbib-style \citep/\citet.
%  Drop into the thesis with \input{thesis_background}.
% ============================================================================

\chapter{Background}
\label{ch:background}

This chapter develops the theory required to follow the rest of the thesis. We
proceed from the general framework of sequential decision making under
uncertainty (\S\ref{sec:bg-rl}--\S\ref{sec:bg-pomdp}), through the specific
learning algorithms used to train our agents (\S\ref{sec:bg-ppo}) and the reward
machinery that shapes them (\S\ref{sec:bg-shaping}), to the model-based world
models that motivate the belief-centric view (\S\ref{sec:bg-worldmodel}).
We then introduce the interpretability tools that this work rests on: the
representational hypotheses that make internal states legible
(\S\ref{sec:bg-representations}), the family of activation-steering interventions
we study (\S\ref{sec:bg-steering}), and finally the belief/behaviour distinction
that frames our central question (\S\ref{sec:bg-belief}).

% ---------------------------------------------------------------------------
\section{Reinforcement Learning and Markov Decision Processes}
\label{sec:bg-rl}

Reinforcement learning (RL) formalises an agent that learns to act by trial and
error so as to maximise a scalar reward \citep{sutton2018reinforcement}. The
standard model is the \emph{Markov Decision Process} (MDP), a tuple
$\mathcal{M} = (\mathcal{S}, \mathcal{A}, P, R, \gamma)$ with state space
$\mathcal{S}$, action space $\mathcal{A}$, transition kernel
$P(s' \mid s, a)$, reward function $R(s,a)$, and discount factor
$\gamma \in [0,1)$. At each step $t$ the agent observes state $s_t$, samples an
action $a_t \sim \pi(\cdot \mid s_t)$ from its \emph{policy} $\pi$, receives
reward $r_t = R(s_t, a_t)$, and transitions to $s_{t+1} \sim P(\cdot \mid s_t,
a_t)$. The objective is to maximise the expected discounted return

\begin{equation}
    G_t \;=\; \sum_{k=0}^{\infty} \gamma^{k}\, r_{t+k},
    \qquad
    J(\pi) \;=\; \mathbb{E}_{\pi}\!\left[G_0\right].
    \label{eq:return}
\end{equation}

The discount $\gamma$ trades off immediate against future reward and keeps the
sum finite. Two value functions summarise long-run desirability: the
\emph{state value} $V^{\pi}(s) = \mathbb{E}_{\pi}[G_t \mid s_t = s]$ and the
\emph{action value} $Q^{\pi}(s,a) = \mathbb{E}_{\pi}[G_t \mid s_t = s, a_t = a]$.
They satisfy the Bellman consistency equations, e.g.

\begin{equation}
    V^{\pi}(s) \;=\; \mathbb{E}_{a\sim\pi,\, s'\sim P}
        \big[\, R(s,a) + \gamma\, V^{\pi}(s') \,\big].
    \label{eq:bellman}
\end{equation}

The \emph{advantage} $A^{\pi}(s,a) = Q^{\pi}(s,a) - V^{\pi}(s)$ measures how much
better action $a$ is than the policy's average behaviour in $s$. For a
parametric policy $\pi_\theta$, the policy-gradient theorem gives an unbiased
ascent direction on $J$:

\begin{equation}
    \nabla_\theta J(\pi_\theta) \;=\;
        \mathbb{E}_{\pi_\theta}\!\left[
            \nabla_\theta \log \pi_\theta(a_t \mid s_t)\; A^{\pi_\theta}(s_t, a_t)
        \right].
    \label{eq:pg}
\end{equation}

Equation~\eqref{eq:pg} is the basis of the actor--critic methods used in this
thesis: a \emph{critic} estimates the advantage and an \emph{actor}
$\pi_\theta$ is nudged to increase the log-probability of advantageous actions.

% ---------------------------------------------------------------------------
\section{Partial Observability and Belief States}
\label{sec:bg-pomdp}

In most realistic environments the agent does not see the full state. This is
captured by the \emph{Partially Observable MDP} (POMDP)
$(\mathcal{S}, \mathcal{A}, P, R, \Omega, O, \gamma)$, which augments the MDP
with an observation space $\Omega$ and an observation model $O(o \mid s)$
\citep{kaelbling1998planning}. The agent receives only $o_t \sim O(\cdot \mid
s_t)$, which need not identify $s_t$. A policy that conditions on the current
observation alone is therefore insufficient; optimal behaviour requires
integrating the \emph{history} $\tau_t = (o_1, a_1, \dots, a_{t-1}, o_t)$.

The central object is the \emph{belief state}, the posterior over hidden states
given the history,

\begin{equation}
    b_t(s) \;=\; \Pr\!\left(s_t = s \,\middle|\, \tau_t\right),
    \label{eq:belief}
\end{equation}

updated recursively by Bayesian filtering,

\begin{equation}
    b_{t}(s') \;\propto\;
        O(o_{t} \mid s')
        \sum_{s \in \mathcal{S}} P(s' \mid s, a_{t-1})\, b_{t-1}(s).
    \label{eq:belief-update}
\end{equation}

The belief $b_t$ is a \emph{sufficient statistic} of the history for predicting
the future: it is the minimal quantity from which the reward-relevant
distribution over states can be recovered \citep{astrom1965optimal}. Crucially,
a POMDP over $\mathcal{S}$ is equivalent to a fully observable
\emph{belief MDP} over the continuous belief simplex $\Delta(\mathcal{S})$, so an
optimal Markov policy $\pi^{\star}(a \mid b)$ exists as a function of the belief
alone. This equivalence is the conceptual backbone of the thesis: \emph{the
belief is the internal variable that must be inferred, maintained, and acted
upon}, and everything we later probe or steer is an approximation to it.

Exact belief tracking is intractable for large $\mathcal{S}$. In deep RL the
belief is approximated by a learned recurrent summary of the history. A
recurrent network with hidden state $h_t \in \mathbb{R}^{d}$ implements

\begin{equation}
    h_t \;=\; f_\theta\!\left(h_{t-1},\, o_t,\, a_{t-1}\right),
    \qquad
    \pi_\theta(a_t \mid \tau_t) \;=\; \pi_\theta(a_t \mid h_t),
    \label{eq:recurrent}
\end{equation}

with $h_t$ trained end-to-end to carry exactly the history information the policy
and value heads need. We use the Gated Recurrent Unit (GRU) for $f_\theta$
\citep{cho2014learning}, whose gating stabilises credit assignment over long
horizons. Under this view $h_t$ is a \emph{learned, distributed belief}: it is
not guaranteed to equal the Bayesian posterior~\eqref{eq:belief}, but a
well-trained agent must encode within $h_t$ whatever latent structure its optimal
policy depends on. Recovering that structure from $h_t$ is the subject of
\S\ref{sec:bg-representations}.

% ---------------------------------------------------------------------------
\section{Proximal Policy Optimization}
\label{sec:bg-ppo}

We train recurrent policies with Proximal Policy Optimization (PPO)
\citep{schulman2017proximal}, a first-order policy-gradient method that is stable
and sample-efficient enough for POMDP navigation. PPO improves on the raw
gradient~\eqref{eq:pg} in two ways.

First, advantages are estimated by \emph{Generalized Advantage Estimation}
(GAE), an exponentially-weighted average of $n$-step temporal-difference errors
that trades bias for variance via $\lambda \in [0,1]$
\citep{schulman2016high}:

\begin{equation}
    \hat{A}_t \;=\; \sum_{l=0}^{\infty} (\gamma\lambda)^{l}\, \delta_{t+l},
    \qquad
    \delta_t \;=\; r_t + \gamma\, V_\phi(s_{t+1}) - V_\phi(s_t),
    \label{eq:gae}
\end{equation}

where $V_\phi$ is the learned critic. Second, PPO constrains each update so the
new policy stays close to the data-collecting policy $\pi_{\theta_{\text{old}}}$,
which prevents destructively large steps. Writing the importance ratio
$\rho_t(\theta) = \pi_\theta(a_t \mid s_t) / \pi_{\theta_{\text{old}}}(a_t \mid
s_t)$, PPO maximises the clipped surrogate

\begin{equation}
    L^{\text{CLIP}}(\theta) \;=\;
        \mathbb{E}_t\!\left[
            \min\!\Big(
                \rho_t(\theta)\, \hat{A}_t,\;
                \operatorname{clip}\!\big(\rho_t(\theta),\, 1-\epsilon,\, 1+\epsilon\big)\, \hat{A}_t
            \Big)
        \right].
    \label{eq:ppo-clip}
\end{equation}

The clip range $\epsilon$ caps the incentive to move $\rho_t$ far from $1$. The
full objective adds a value-regression term and an entropy bonus,

\begin{equation}
    L(\theta,\phi) \;=\;
        L^{\text{CLIP}}(\theta)
        \;-\; c_v\, \mathbb{E}_t\!\big[(V_\phi(s_t) - \hat{G}_t)^2\big]
        \;+\; c_e\, \mathbb{E}_t\!\big[\mathcal{H}[\pi_\theta(\cdot \mid s_t)]\big],
    \label{eq:ppo-full}
\end{equation}

where $\hat{G}_t = \hat{A}_t + V_\phi(s_t)$ is the value target and
$\mathcal{H}$ denotes entropy. The entropy coefficient $c_e$ controls
exploration; as we show in the experiments, it is decisive when the task admits
a lazy, belief-free local optimum, because escaping that optimum requires
sustained exploration of belief-conditional behaviour. For recurrent policies,
equations~\eqref{eq:gae}--\eqref{eq:ppo-full} are optimised over sequences: the
hidden state~\eqref{eq:recurrent} is carried through each rollout and reset at
episode boundaries, so gradients flow through time via backpropagation through
the unrolled GRU.

% ---------------------------------------------------------------------------
\section{Reward Shaping and Potential-Based Invariance}
\label{sec:bg-shaping}

Sparse rewards make credit assignment hard: if reward arrives only at a distant
goal, the gradient~\eqref{eq:pg} is uninformative for most of an episode.
\emph{Reward shaping} adds an auxiliary signal $F(s,a,s')$ to densify learning.
The danger is that a poorly chosen $F$ changes what is optimal, so the agent
learns to exploit the shaping instead of solving the task. \citet{ng1999policy}
show that this is avoided \emph{if and only if} the shaping is
\emph{potential-based}: for some potential function $\Phi : \mathcal{S} \to
\mathbb{R}$,

\begin{equation}
    F(s,a,s') \;=\; \gamma\, \Phi(s') - \Phi(s).
    \label{eq:pbrs}
\end{equation}

Under \eqref{eq:pbrs} the set of optimal policies is provably unchanged, because
the added term telescopes: over any trajectory the shaping contributes only
$\gamma^{T}\Phi(s_T) - \Phi(s_0)$, a quantity that depends on the endpoints, not
on the path taken. Potential-based reward shaping (PBRS) therefore accelerates
learning without biasing the solution.

In this thesis the potential is derived from a \emph{cost-to-go} oracle. Let
$\mathrm{ctg}(s)$ be the minimum number of primitive actions needed to reach the
goal from $s$, computed by a shortest-path search on the (known) map; we set

\begin{equation}
    \Phi(s) \;=\; -\,\mathrm{ctg}(s).
    \label{eq:phi-ctg}
\end{equation}

Moving closer to the goal ($\mathrm{ctg}$ decreasing) yields positive shaping,
providing a dense navigational gradient. Two properties of PBRS matter for our
analysis. First, because the shaping is invariant, it \emph{cannot} manufacture
a preference the terminal reward does not already encode; it can only make an
existing preference easier to discover. Second, its contribution to the
\emph{return difference} between two candidate endpoints is bounded by
$\lvert \Phi(s_A) - \Phi(s_B)\rvert$, which for a small
$\mathrm{ctg}$ gap can be negligible relative to a large terminal bonus. We
return to both points when explaining why shaping toward a target does not, on
its own, force a partially observable agent to condition its behaviour on the
correct latent belief.

% ---------------------------------------------------------------------------
\section{World Models and the Recurrent State-Space Model}
\label{sec:bg-worldmodel}

An alternative to learning $h_t$ purely from the RL loss is to learn an explicit
generative \emph{world model} that predicts observations and rewards, and to
train the policy inside the model's imagined rollouts. DreamerV3
\citep{hafner2023mastering,hafner2020dream} is the state-of-the-art instance and
underlies the model-based agents in this work. Its core is the \emph{Recurrent
State-Space Model} (RSSM), which factorises the latent state into a
deterministic recurrent path $h_t$ and a stochastic latent $z_t$:

\begin{align}
    h_t         &= f_\theta(h_{t-1}, z_{t-1}, a_{t-1}) && \text{(recurrent model)} \label{eq:rssm-h}\\
    \hat{z}_t   &\sim p_\theta(\hat{z}_t \mid h_t)     && \text{(prior / transition)} \label{eq:rssm-prior}\\
    z_t         &\sim q_\theta(z_t \mid h_t, x_t)      && \text{(posterior / filter)} \label{eq:rssm-post}
\end{align}

where $x_t$ is the encoded observation. The \emph{prior}
$p_\theta$~\eqref{eq:rssm-prior} predicts the next latent from dynamics alone
(enabling imagination), while the \emph{posterior}
$q_\theta$~\eqref{eq:rssm-post} corrects it using the actual observation --- an
amortised, learned version of the Bayesian filter~\eqref{eq:belief-update}. The
joint latent $s_t = (h_t, z_t)$ is the model's belief state. Decoder heads
reconstruct the observation $\hat{x}_t$ and predict reward from $s_t$, and the
whole model is trained by maximising an evidence lower bound of the form

\begin{equation}
    \mathcal{L}(\theta) \;=\;
        \underbrace{\mathbb{E}_{q}\big[\log p_\theta(x_t \mid s_t)\big]}_{\text{reconstruction}}
        \;-\; \beta\,
        \underbrace{\mathrm{KL}\!\big[\,q_\theta(z_t \mid h_t, x_t)\,\big\|\,p_\theta(\hat z_t \mid h_t)\,\big]}_{\text{dynamics / representation}}.
    \label{eq:elbo}
\end{equation}

The actor and critic are trained on trajectories \emph{imagined} by rolling the
prior forward without observations, then decoded back to interpretable
predictions. Two features of this architecture are central to the thesis.
\textbf{(i)} The belief $(h_t, z_t)$ is a single shared object consumed by
multiple heads --- the actor, the reward predictor, and the observation decoder
--- so there is no architecturally separate ``skill'' subspace. \textbf{(ii)}
Because the decoder maps the latent back to observation space, the model exposes
an explicit \emph{belief readout}: one can decode what the agent ``believes'' the
world looks like directly from $s_t$, and watch how that decoded belief changes
under intervention. This makes the world model an unusually transparent
substrate for the interpretability questions we now formalise.

% ---------------------------------------------------------------------------
\section{Representations, Probing, and the Linear Hypothesis}
\label{sec:bg-representations}

Neural agents encode task-relevant variables in \emph{distributed}
representations: a concept is not a single unit but a pattern across many units
\citep{rumelhart1986learning}. The empirically influential \emph{linear
representation hypothesis} (LRH) sharpens this: many high-level concepts are
encoded as \emph{directions} in activation space, so that the presence or value
of a concept is (approximately) an affine function of the hidden state
\citep{mikolov2013linguistic,park2024linear}. Under the LRH, a hidden state
$h \in \mathbb{R}^{d}$ decomposes as a sum of concept directions plus residual,
and the value of concept $i$ is read off by projecting onto a direction $v_i$.

The standard tool for testing this is a \emph{linear probe}: a simple classifier
trained to predict a labelled latent variable $y$ from a frozen activation $h$
\citep{alain2017understanding}. For a categorical belief label with classes
$c$, the probe is a linear--softmax map

\begin{equation}
    \hat{p}(y = c \mid h) \;=\; \operatorname{softmax}_c\!\big(W h + b\big),
    \label{eq:probe}
\end{equation}

and high held-out accuracy is taken as evidence that the information $y$ is
linearly \emph{decodable} from $h$. A key methodological caution, which this
thesis makes concrete, is that decodability is a statement about what a
\emph{downstream reader} could in principle extract --- it does \emph{not} imply
that the agent's own policy head uses that direction. In our agents we attach
such a probe as an auxiliary head trained jointly with the policy (a
``belief head''); it provides a free, continuously-monitored readout of the
latent belief without feeding back into the actor, so that decodability and use
can be measured separately. When a concept is linearly represented, a canonical
estimate of its direction is the \emph{difference of class means} in activation
space,

\begin{equation}
    v \;=\; \mu_{+} - \mu_{-}, \qquad
    \mu_{\pm} \;=\; \frac{1}{|\mathcal{D}_{\pm}|} \sum_{h \in \mathcal{D}_{\pm}} h,
    \label{eq:diffmeans}
\end{equation}

where $\mathcal{D}_{+}$ and $\mathcal{D}_{-}$ are activations collected under the
two values of the concept \citep{marks2024geometry}. The vector $v$ points from
one class centroid to the other and is the workhorse direction for the
interventions in the next section. The LRH is, however, only approximate:
\emph{superposition} lets networks pack more features than dimensions by
overlapping them in non-orthogonal directions \citep{elhage2022superposition},
and some features occupy multi-dimensional subspaces rather than single lines.
This non-orthogonality is precisely what makes targeted intervention leaky, a
theme we develop in \S\ref{sec:bg-belief}.

% ---------------------------------------------------------------------------
\section{Activation Steering}
\label{sec:bg-steering}

\emph{Activation steering} intervenes on a model's computation at inference time
by editing its hidden states along a chosen direction, rather than retraining
weights. Given a steering direction $v$ (often estimated by
\eqref{eq:diffmeans}), a layer activation $h$, and a scale $\alpha$, the
canonical edit is additive:

\begin{equation}
    h' \;=\; h + \alpha\, v.
    \label{eq:steer-add}
\end{equation}

This is \emph{activation addition} (ActAdd), which injects a contrast vector at
inference to push behaviour toward or away from a target concept
\citep{turner2023activation}. \emph{Contrastive Activation Addition} (CAA)
constructs $v$ as the mean difference of activations over matched positive and
negative prompt pairs \citep{rimsky2024steering}, i.e.\ a structured instance of
\eqref{eq:diffmeans}. A complementary operation is directional \emph{ablation},
which removes a concept by projecting it out of the activation,

\begin{equation}
    h' \;=\; h - (\hat{v}^{\top} h)\, \hat{v}, \qquad \hat{v} = v / \lVert v \rVert,
    \label{eq:steer-ablate}
\end{equation}

as used to suppress the refusal direction in language models
\citep{arditi2024refusal}. Finer-grained variants steer individual features
discovered by \emph{sparse autoencoders}, which decompose $h$ into a sparse,
ideally monosemantic dictionary and amplify or suppress single entries
\citep{bricken2023monosemanticity,cunningham2024sparse}. In this thesis we apply
\eqref{eq:steer-add}--\eqref{eq:steer-ablate} to the recurrent belief carrier
$h_t$ of an RL agent: the direction $v$ is a difference-of-means over the latent
belief (e.g.\ activations on maps of one category minus another), and it is
injected over a window of timesteps while the agent acts. The behavioural
consequence --- whether the intervention changes the agent's downstream choice
--- is then read out from its actions, and the belief consequence from the
decoded world model or belief probe.

Every method in this family shares an implicit assumption: that editing along
$v$ leaves the computations served by other directions untouched. Because
representations are entangled and non-orthogonal
(\S\ref{sec:bg-representations}), any edit with a nonzero projection onto other
functional directions perturbs them too. Whether a given belief direction is
in fact \emph{used} by the policy, and whether steering it produces a clean
behavioural change or leaks into unrelated capabilities, are empirical questions
--- and the ones this thesis is built to answer.

% ---------------------------------------------------------------------------
\section{Belief, Skill, and the Decoupling Question}
\label{sec:bg-belief}

We can now state the conceptual frame precisely. We distinguish two readouts of a
shared latent state $z$ (either the GRU state $h_t$ or the RSSM latent
$(h_t,z_t)$). The \emph{belief readout} $B(z)$ is the model's inferred
distribution over hidden world structure --- e.g.\ the decoded observation of a
world model, or the belief probe~\eqref{eq:probe}. The \emph{policy readout}
$\pi(z)$ is the distribution over actions. A latent variable is
\emph{behaviourally load-bearing} when the policy's optimal action depends on it;
it is \emph{merely decodable} when $B(z)$ recovers it but $\pi(z)$ does not use
it.

Localised steering implicitly assumes these are separable: that one can edit the
belief without disturbing the policy, or edit the policy's target behaviour
without disturbing the belief. Following the diagnostic used in this work, an
intervention $\delta$ applied to $z$ is \emph{belief-preserving} iff it moves the
policy while leaving the belief readout intact:

\begin{equation}
    \big\lVert B(z + \delta) - B(z) \big\rVert < \varepsilon
    \quad\text{and}\quad
    \big\lVert \pi(z + \delta) - \pi(z) \big\rVert > \eta,
    \label{eq:belief-invariance}
\end{equation}

for pre-specified tolerances $\varepsilon$ (permitted belief drift) and $\eta$
(required behavioural change). Symmetrically, a behaviour-targeting edit is
\emph{clean} only if it does not silently shift $B(z)$. When belief and policy
readouts load onto overlapping directions of a shared latent, no strictly linear
edit can satisfy \eqref{eq:belief-invariance} exactly; the severity of the
resulting cross-talk is a property of the representational geometry, not of the
method. This is why architectures with an explicit belief decoder are valuable:
they let $B(z)$ be measured directly, turning \eqref{eq:belief-invariance} from a
philosophical distinction into a checkable quantity.

The environments and agents in the following chapters are designed to make this
question concrete and controllable. The agent must infer a latent map property
(the belief) from a partial history and select among behaviours accordingly (the
skill), while an auxiliary belief head~\eqref{eq:probe} and, for the world-model
agents, an observation decoder provide independent readouts of $B(z)$. This
setup lets us ask, and answer empirically, the two questions that motivate the
thesis: \emph{(i)} when is a decodable belief actually used by the policy, and
\emph{(ii)} does steering the belief direction change behaviour as intended, or
does it drift the belief while leaving behaviour untouched --- or vice versa.
.
% ============================================================================

\chapter{Background}
\label{ch:background}

This chapter develops the theory required to follow the rest of the thesis. We
proceed from the general framework of sequential decision making under
uncertainty (\S\ref{sec:bg-rl}--\S\ref{sec:bg-pomdp}), through the specific
learning algorithms used to train our agents (\S\ref{sec:bg-ppo}) and the reward
machinery that shapes them (\S\ref{sec:bg-shaping}), to the model-based world
models that motivate the belief-centric view (\S\ref{sec:bg-worldmodel}).
We then introduce the interpretability tools that this work rests on: the
representational hypotheses that make internal states legible
(\S\ref{sec:bg-representations}), the family of activation-steering interventions
we study (\S\ref{sec:bg-steering}), and finally the belief/behaviour distinction
that frames our central question (\S\ref{sec:bg-belief}).

% ---------------------------------------------------------------------------
\section{Reinforcement Learning and Markov Decision Processes}
\label{sec:bg-rl}

Reinforcement learning (RL) formalises an agent that learns to act by trial and
error so as to maximise a scalar reward \citep{sutton2018reinforcement}. The
standard model is the \emph{Markov Decision Process} (MDP), a tuple
$\mathcal{M} = (\mathcal{S}, \mathcal{A}, P, R, \gamma)$ with state space
$\mathcal{S}$, action space $\mathcal{A}$, transition kernel
$P(s' \mid s, a)$, reward function $R(s,a)$, and discount factor
$\gamma \in [0,1)$. At each step $t$ the agent observes state $s_t$, samples an
action $a_t \sim \pi(\cdot \mid s_t)$ from its \emph{policy} $\pi$, receives
reward $r_t = R(s_t, a_t)$, and transitions to $s_{t+1} \sim P(\cdot \mid s_t,
a_t)$. The objective is to maximise the expected discounted return

\begin{equation}
    G_t \;=\; \sum_{k=0}^{\infty} \gamma^{k}\, r_{t+k},
    \qquad
    J(\pi) \;=\; \mathbb{E}_{\pi}\!\left[G_0\right].
    \label{eq:return}
\end{equation}

The discount $\gamma$ trades off immediate against future reward and keeps the
sum finite. Two value functions summarise long-run desirability: the
\emph{state value} $V^{\pi}(s) = \mathbb{E}_{\pi}[G_t \mid s_t = s]$ and the
\emph{action value} $Q^{\pi}(s,a) = \mathbb{E}_{\pi}[G_t \mid s_t = s, a_t = a]$.
They satisfy the Bellman consistency equations, e.g.

\begin{equation}
    V^{\pi}(s) \;=\; \mathbb{E}_{a\sim\pi,\, s'\sim P}
        \big[\, R(s,a) + \gamma\, V^{\pi}(s') \,\big].
    \label{eq:bellman}
\end{equation}

The \emph{advantage} $A^{\pi}(s,a) = Q^{\pi}(s,a) - V^{\pi}(s)$ measures how much
better action $a$ is than the policy's average behaviour in $s$. For a
parametric policy $\pi_\theta$, the policy-gradient theorem gives an unbiased
ascent direction on $J$:

\begin{equation}
    \nabla_\theta J(\pi_\theta) \;=\;
        \mathbb{E}_{\pi_\theta}\!\left[
            \nabla_\theta \log \pi_\theta(a_t \mid s_t)\; A^{\pi_\theta}(s_t, a_t)
        \right].
    \label{eq:pg}
\end{equation}

Equation~\eqref{eq:pg} is the basis of the actor--critic methods used in this
thesis: a \emph{critic} estimates the advantage and an \emph{actor}
$\pi_\theta$ is nudged to increase the log-probability of advantageous actions.

% ---------------------------------------------------------------------------
\section{Partial Observability and Belief States}
\label{sec:bg-pomdp}

In most realistic environments the agent does not see the full state. This is
captured by the \emph{Partially Observable MDP} (POMDP)
$(\mathcal{S}, \mathcal{A}, P, R, \Omega, O, \gamma)$, which augments the MDP
with an observation space $\Omega$ and an observation model $O(o \mid s)$
\citep{kaelbling1998planning}. The agent receives only $o_t \sim O(\cdot \mid
s_t)$, which need not identify $s_t$. A policy that conditions on the current
observation alone is therefore insufficient; optimal behaviour requires
integrating the \emph{history} $\tau_t = (o_1, a_1, \dots, a_{t-1}, o_t)$.

The central object is the \emph{belief state}, the posterior over hidden states
given the history,

\begin{equation}
    b_t(s) \;=\; \Pr\!\left(s_t = s \,\middle|\, \tau_t\right),
    \label{eq:belief}
\end{equation}

updated recursively by Bayesian filtering,

\begin{equation}
    b_{t}(s') \;\propto\;
        O(o_{t} \mid s')
        \sum_{s \in \mathcal{S}} P(s' \mid s, a_{t-1})\, b_{t-1}(s).
    \label{eq:belief-update}
\end{equation}

The belief $b_t$ is a \emph{sufficient statistic} of the history for predicting
the future: it is the minimal quantity from which the reward-relevant
distribution over states can be recovered \citep{astrom1965optimal}. Crucially,
a POMDP over $\mathcal{S}$ is equivalent to a fully observable
\emph{belief MDP} over the continuous belief simplex $\Delta(\mathcal{S})$, so an
optimal Markov policy $\pi^{\star}(a \mid b)$ exists as a function of the belief
alone. This equivalence is the conceptual backbone of the thesis: \emph{the
belief is the internal variable that must be inferred, maintained, and acted
upon}, and everything we later probe or steer is an approximation to it.

Exact belief tracking is intractable for large $\mathcal{S}$. In deep RL the
belief is approximated by a learned recurrent summary of the history. A
recurrent network with hidden state $h_t \in \mathbb{R}^{d}$ implements

\begin{equation}
    h_t \;=\; f_\theta\!\left(h_{t-1},\, o_t,\, a_{t-1}\right),
    \qquad
    \pi_\theta(a_t \mid \tau_t) \;=\; \pi_\theta(a_t \mid h_t),
    \label{eq:recurrent}
\end{equation}

with $h_t$ trained end-to-end to carry exactly the history information the policy
and value heads need. We use the Gated Recurrent Unit (GRU) for $f_\theta$
\citep{cho2014learning}, whose gating stabilises credit assignment over long
horizons. Under this view $h_t$ is a \emph{learned, distributed belief}: it is
not guaranteed to equal the Bayesian posterior~\eqref{eq:belief}, but a
well-trained agent must encode within $h_t$ whatever latent structure its optimal
policy depends on. Recovering that structure from $h_t$ is the subject of
\S\ref{sec:bg-representations}.

% ---------------------------------------------------------------------------
\section{Proximal Policy Optimization}
\label{sec:bg-ppo}

We train recurrent policies with Proximal Policy Optimization (PPO)
\citep{schulman2017proximal}, a first-order policy-gradient method that is stable
and sample-efficient enough for POMDP navigation. PPO improves on the raw
gradient~\eqref{eq:pg} in two ways.

First, advantages are estimated by \emph{Generalized Advantage Estimation}
(GAE), an exponentially-weighted average of $n$-step temporal-difference errors
that trades bias for variance via $\lambda \in [0,1]$
\citep{schulman2016high}:

\begin{equation}
    \hat{A}_t \;=\; \sum_{l=0}^{\infty} (\gamma\lambda)^{l}\, \delta_{t+l},
    \qquad
    \delta_t \;=\; r_t + \gamma\, V_\phi(s_{t+1}) - V_\phi(s_t),
    \label{eq:gae}
\end{equation}

where $V_\phi$ is the learned critic. Second, PPO constrains each update so the
new policy stays close to the data-collecting policy $\pi_{\theta_{\text{old}}}$,
which prevents destructively large steps. Writing the importance ratio
$\rho_t(\theta) = \pi_\theta(a_t \mid s_t) / \pi_{\theta_{\text{old}}}(a_t \mid
s_t)$, PPO maximises the clipped surrogate

\begin{equation}
    L^{\text{CLIP}}(\theta) \;=\;
        \mathbb{E}_t\!\left[
            \min\!\Big(
                \rho_t(\theta)\, \hat{A}_t,\;
                \operatorname{clip}\!\big(\rho_t(\theta),\, 1-\epsilon,\, 1+\epsilon\big)\, \hat{A}_t
            \Big)
        \right].
    \label{eq:ppo-clip}
\end{equation}

The clip range $\epsilon$ caps the incentive to move $\rho_t$ far from $1$. The
full objective adds a value-regression term and an entropy bonus,

\begin{equation}
    L(\theta,\phi) \;=\;
        L^{\text{CLIP}}(\theta)
        \;-\; c_v\, \mathbb{E}_t\!\big[(V_\phi(s_t) - \hat{G}_t)^2\big]
        \;+\; c_e\, \mathbb{E}_t\!\big[\mathcal{H}[\pi_\theta(\cdot \mid s_t)]\big],
    \label{eq:ppo-full}
\end{equation}

where $\hat{G}_t = \hat{A}_t + V_\phi(s_t)$ is the value target and
$\mathcal{H}$ denotes entropy. The entropy coefficient $c_e$ controls
exploration; as we show in the experiments, it is decisive when the task admits
a lazy, belief-free local optimum, because escaping that optimum requires
sustained exploration of belief-conditional behaviour. For recurrent policies,
equations~\eqref{eq:gae}--\eqref{eq:ppo-full} are optimised over sequences: the
hidden state~\eqref{eq:recurrent} is carried through each rollout and reset at
episode boundaries, so gradients flow through time via backpropagation through
the unrolled GRU.

% ---------------------------------------------------------------------------
\section{Reward Shaping and Potential-Based Invariance}
\label{sec:bg-shaping}

Sparse rewards make credit assignment hard: if reward arrives only at a distant
goal, the gradient~\eqref{eq:pg} is uninformative for most of an episode.
\emph{Reward shaping} adds an auxiliary signal $F(s,a,s')$ to densify learning.
The danger is that a poorly chosen $F$ changes what is optimal, so the agent
learns to exploit the shaping instead of solving the task. \citet{ng1999policy}
show that this is avoided \emph{if and only if} the shaping is
\emph{potential-based}: for some potential function $\Phi : \mathcal{S} \to
\mathbb{R}$,

\begin{equation}
    F(s,a,s') \;=\; \gamma\, \Phi(s') - \Phi(s).
    \label{eq:pbrs}
\end{equation}

Under \eqref{eq:pbrs} the set of optimal policies is provably unchanged, because
the added term telescopes: over any trajectory the shaping contributes only
$\gamma^{T}\Phi(s_T) - \Phi(s_0)$, a quantity that depends on the endpoints, not
on the path taken. Potential-based reward shaping (PBRS) therefore accelerates
learning without biasing the solution.

In this thesis the potential is derived from a \emph{cost-to-go} oracle. Let
$\mathrm{ctg}(s)$ be the minimum number of primitive actions needed to reach the
goal from $s$, computed by a shortest-path search on the (known) map; we set

\begin{equation}
    \Phi(s) \;=\; -\,\mathrm{ctg}(s).
    \label{eq:phi-ctg}
\end{equation}

Moving closer to the goal ($\mathrm{ctg}$ decreasing) yields positive shaping,
providing a dense navigational gradient. Two properties of PBRS matter for our
analysis. First, because the shaping is invariant, it \emph{cannot} manufacture
a preference the terminal reward does not already encode; it can only make an
existing preference easier to discover. Second, its contribution to the
\emph{return difference} between two candidate endpoints is bounded by
$\lvert \Phi(s_A) - \Phi(s_B)\rvert$, which for a small
$\mathrm{ctg}$ gap can be negligible relative to a large terminal bonus. We
return to both points when explaining why shaping toward a target does not, on
its own, force a partially observable agent to condition its behaviour on the
correct latent belief.

% ---------------------------------------------------------------------------
\section{World Models and the Recurrent State-Space Model}
\label{sec:bg-worldmodel}

An alternative to learning $h_t$ purely from the RL loss is to learn an explicit
generative \emph{world model} that predicts observations and rewards, and to
train the policy inside the model's imagined rollouts. DreamerV3
\citep{hafner2023mastering,hafner2020dream} is the state-of-the-art instance and
underlies the model-based agents in this work. Its core is the \emph{Recurrent
State-Space Model} (RSSM), which factorises the latent state into a
deterministic recurrent path $h_t$ and a stochastic latent $z_t$:

\begin{align}
    h_t         &= f_\theta(h_{t-1}, z_{t-1}, a_{t-1}) && \text{(recurrent model)} \label{eq:rssm-h}\\
    \hat{z}_t   &\sim p_\theta(\hat{z}_t \mid h_t)     && \text{(prior / transition)} \label{eq:rssm-prior}\\
    z_t         &\sim q_\theta(z_t \mid h_t, x_t)      && \text{(posterior / filter)} \label{eq:rssm-post}
\end{align}

where $x_t$ is the encoded observation. The \emph{prior}
$p_\theta$~\eqref{eq:rssm-prior} predicts the next latent from dynamics alone
(enabling imagination), while the \emph{posterior}
$q_\theta$~\eqref{eq:rssm-post} corrects it using the actual observation --- an
amortised, learned version of the Bayesian filter~\eqref{eq:belief-update}. The
joint latent $s_t = (h_t, z_t)$ is the model's belief state. Decoder heads
reconstruct the observation $\hat{x}_t$ and predict reward from $s_t$, and the
whole model is trained by maximising an evidence lower bound of the form

\begin{equation}
    \mathcal{L}(\theta) \;=\;
        \underbrace{\mathbb{E}_{q}\big[\log p_\theta(x_t \mid s_t)\big]}_{\text{reconstruction}}
        \;-\; \beta\,
        \underbrace{\mathrm{KL}\!\big[\,q_\theta(z_t \mid h_t, x_t)\,\big\|\,p_\theta(\hat z_t \mid h_t)\,\big]}_{\text{dynamics / representation}}.
    \label{eq:elbo}
\end{equation}

The actor and critic are trained on trajectories \emph{imagined} by rolling the
prior forward without observations, then decoded back to interpretable
predictions. Two features of this architecture are central to the thesis.
\textbf{(i)} The belief $(h_t, z_t)$ is a single shared object consumed by
multiple heads --- the actor, the reward predictor, and the observation decoder
--- so there is no architecturally separate ``skill'' subspace. \textbf{(ii)}
Because the decoder maps the latent back to observation space, the model exposes
an explicit \emph{belief readout}: one can decode what the agent ``believes'' the
world looks like directly from $s_t$, and watch how that decoded belief changes
under intervention. This makes the world model an unusually transparent
substrate for the interpretability questions we now formalise.

% ---------------------------------------------------------------------------
\section{Representations, Probing, and the Linear Hypothesis}
\label{sec:bg-representations}

Neural agents encode task-relevant variables in \emph{distributed}
representations: a concept is not a single unit but a pattern across many units
\citep{rumelhart1986learning}. The empirically influential \emph{linear
representation hypothesis} (LRH) sharpens this: many high-level concepts are
encoded as \emph{directions} in activation space, so that the presence or value
of a concept is (approximately) an affine function of the hidden state
\citep{mikolov2013linguistic,park2024linear}. Under the LRH, a hidden state
$h \in \mathbb{R}^{d}$ decomposes as a sum of concept directions plus residual,
and the value of concept $i$ is read off by projecting onto a direction $v_i$.

The standard tool for testing this is a \emph{linear probe}: a simple classifier
trained to predict a labelled latent variable $y$ from a frozen activation $h$
\citep{alain2017understanding}. For a categorical belief label with classes
$c$, the probe is a linear--softmax map

\begin{equation}
    \hat{p}(y = c \mid h) \;=\; \operatorname{softmax}_c\!\big(W h + b\big),
    \label{eq:probe}
\end{equation}

and high held-out accuracy is taken as evidence that the information $y$ is
linearly \emph{decodable} from $h$. A key methodological caution, which this
thesis makes concrete, is that decodability is a statement about what a
\emph{downstream reader} could in principle extract --- it does \emph{not} imply
that the agent's own policy head uses that direction. In our agents we attach
such a probe as an auxiliary head trained jointly with the policy (a
``belief head''); it provides a free, continuously-monitored readout of the
latent belief without feeding back into the actor, so that decodability and use
can be measured separately. When a concept is linearly represented, a canonical
estimate of its direction is the \emph{difference of class means} in activation
space,

\begin{equation}
    v \;=\; \mu_{+} - \mu_{-}, \qquad
    \mu_{\pm} \;=\; \frac{1}{|\mathcal{D}_{\pm}|} \sum_{h \in \mathcal{D}_{\pm}} h,
    \label{eq:diffmeans}
\end{equation}

where $\mathcal{D}_{+}$ and $\mathcal{D}_{-}$ are activations collected under the
two values of the concept \citep{marks2024geometry}. The vector $v$ points from
one class centroid to the other and is the workhorse direction for the
interventions in the next section. The LRH is, however, only approximate:
\emph{superposition} lets networks pack more features than dimensions by
overlapping them in non-orthogonal directions \citep{elhage2022superposition},
and some features occupy multi-dimensional subspaces rather than single lines.
This non-orthogonality is precisely what makes targeted intervention leaky, a
theme we develop in \S\ref{sec:bg-belief}.

% ---------------------------------------------------------------------------
\section{Activation Steering}
\label{sec:bg-steering}

\emph{Activation steering} intervenes on a model's computation at inference time
by editing its hidden states along a chosen direction, rather than retraining
weights. Given a steering direction $v$ (often estimated by
\eqref{eq:diffmeans}), a layer activation $h$, and a scale $\alpha$, the
canonical edit is additive:

\begin{equation}
    h' \;=\; h + \alpha\, v.
    \label{eq:steer-add}
\end{equation}

This is \emph{activation addition} (ActAdd), which injects a contrast vector at
inference to push behaviour toward or away from a target concept
\citep{turner2023activation}. \emph{Contrastive Activation Addition} (CAA)
constructs $v$ as the mean difference of activations over matched positive and
negative prompt pairs \citep{rimsky2024steering}, i.e.\ a structured instance of
\eqref{eq:diffmeans}. A complementary operation is directional \emph{ablation},
which removes a concept by projecting it out of the activation,

\begin{equation}
    h' \;=\; h - (\hat{v}^{\top} h)\, \hat{v}, \qquad \hat{v} = v / \lVert v \rVert,
    \label{eq:steer-ablate}
\end{equation}

as used to suppress the refusal direction in language models
\citep{arditi2024refusal}. Finer-grained variants steer individual features
discovered by \emph{sparse autoencoders}, which decompose $h$ into a sparse,
ideally monosemantic dictionary and amplify or suppress single entries
\citep{bricken2023monosemanticity,cunningham2024sparse}. In this thesis we apply
\eqref{eq:steer-add}--\eqref{eq:steer-ablate} to the recurrent belief carrier
$h_t$ of an RL agent: the direction $v$ is a difference-of-means over the latent
belief (e.g.\ activations on maps of one category minus another), and it is
injected over a window of timesteps while the agent acts. The behavioural
consequence --- whether the intervention changes the agent's downstream choice
--- is then read out from its actions, and the belief consequence from the
decoded world model or belief probe.

Every method in this family shares an implicit assumption: that editing along
$v$ leaves the computations served by other directions untouched. Because
representations are entangled and non-orthogonal
(\S\ref{sec:bg-representations}), any edit with a nonzero projection onto other
functional directions perturbs them too. Whether a given belief direction is
in fact \emph{used} by the policy, and whether steering it produces a clean
behavioural change or leaks into unrelated capabilities, are empirical questions
--- and the ones this thesis is built to answer.

% ---------------------------------------------------------------------------
\section{Belief, Skill, and the Decoupling Question}
\label{sec:bg-belief}

We can now state the conceptual frame precisely. We distinguish two readouts of a
shared latent state $z$ (either the GRU state $h_t$ or the RSSM latent
$(h_t,z_t)$). The \emph{belief readout} $B(z)$ is the model's inferred
distribution over hidden world structure --- e.g.\ the decoded observation of a
world model, or the belief probe~\eqref{eq:probe}. The \emph{policy readout}
$\pi(z)$ is the distribution over actions. A latent variable is
\emph{behaviourally load-bearing} when the policy's optimal action depends on it;
it is \emph{merely decodable} when $B(z)$ recovers it but $\pi(z)$ does not use
it.

Localised steering implicitly assumes these are separable: that one can edit the
belief without disturbing the policy, or edit the policy's target behaviour
without disturbing the belief. Following the diagnostic used in this work, an
intervention $\delta$ applied to $z$ is \emph{belief-preserving} iff it moves the
policy while leaving the belief readout intact:

\begin{equation}
    \big\lVert B(z + \delta) - B(z) \big\rVert < \varepsilon
    \quad\text{and}\quad
    \big\lVert \pi(z + \delta) - \pi(z) \big\rVert > \eta,
    \label{eq:belief-invariance}
\end{equation}

for pre-specified tolerances $\varepsilon$ (permitted belief drift) and $\eta$
(required behavioural change). Symmetrically, a behaviour-targeting edit is
\emph{clean} only if it does not silently shift $B(z)$. When belief and policy
readouts load onto overlapping directions of a shared latent, no strictly linear
edit can satisfy \eqref{eq:belief-invariance} exactly; the severity of the
resulting cross-talk is a property of the representational geometry, not of the
method. This is why architectures with an explicit belief decoder are valuable:
they let $B(z)$ be measured directly, turning \eqref{eq:belief-invariance} from a
philosophical distinction into a checkable quantity.

The environments and agents in the following chapters are designed to make this
question concrete and controllable. The agent must infer a latent map property
(the belief) from a partial history and select among behaviours accordingly (the
skill), while an auxiliary belief head~\eqref{eq:probe} and, for the world-model
agents, an observation decoder provide independent readouts of $B(z)$. This
setup lets us ask, and answer empirically, the two questions that motivate the
thesis: \emph{(i)} when is a decodable belief actually used by the policy, and
\emph{(ii)} does steering the belief direction change behaviour as intended, or
does it drift the belief while leaving behaviour untouched --- or vice versa.
.
% ============================================================================

\chapter{Background}
\label{ch:background}

This chapter develops the theory required to follow the rest of the thesis. We
proceed from the general framework of sequential decision making under
uncertainty (\S\ref{sec:bg-rl}--\S\ref{sec:bg-pomdp}), through the specific
learning algorithms used to train our agents (\S\ref{sec:bg-ppo}) and the reward
machinery that shapes them (\S\ref{sec:bg-shaping}), to the model-based world
models that motivate the belief-centric view (\S\ref{sec:bg-worldmodel}).
We then introduce the interpretability tools that this work rests on: the
representational hypotheses that make internal states legible
(\S\ref{sec:bg-representations}), the family of activation-steering interventions
we study (\S\ref{sec:bg-steering}), and finally the belief/behaviour distinction
that frames our central question (\S\ref{sec:bg-belief}).

% ---------------------------------------------------------------------------
\section{Reinforcement Learning and Markov Decision Processes}
\label{sec:bg-rl}

Reinforcement learning (RL) formalises an agent that learns to act by trial and
error so as to maximise a scalar reward \citep{sutton2018reinforcement}. The
standard model is the \emph{Markov Decision Process} (MDP), a tuple
$\mathcal{M} = (\mathcal{S}, \mathcal{A}, P, R, \gamma)$ with state space
$\mathcal{S}$, action space $\mathcal{A}$, transition kernel
$P(s' \mid s, a)$, reward function $R(s,a)$, and discount factor
$\gamma \in [0,1)$. At each step $t$ the agent observes state $s_t$, samples an
action $a_t \sim \pi(\cdot \mid s_t)$ from its \emph{policy} $\pi$, receives
reward $r_t = R(s_t, a_t)$, and transitions to $s_{t+1} \sim P(\cdot \mid s_t,
a_t)$. The objective is to maximise the expected discounted return

\begin{equation}
    G_t \;=\; \sum_{k=0}^{\infty} \gamma^{k}\, r_{t+k},
    \qquad
    J(\pi) \;=\; \mathbb{E}_{\pi}\!\left[G_0\right].
    \label{eq:return}
\end{equation}

The discount $\gamma$ trades off immediate against future reward and keeps the
sum finite. Two value functions summarise long-run desirability: the
\emph{state value} $V^{\pi}(s) = \mathbb{E}_{\pi}[G_t \mid s_t = s]$ and the
\emph{action value} $Q^{\pi}(s,a) = \mathbb{E}_{\pi}[G_t \mid s_t = s, a_t = a]$.
They satisfy the Bellman consistency equations, e.g.

\begin{equation}
    V^{\pi}(s) \;=\; \mathbb{E}_{a\sim\pi,\, s'\sim P}
        \big[\, R(s,a) + \gamma\, V^{\pi}(s') \,\big].
    \label{eq:bellman}
\end{equation}

The \emph{advantage} $A^{\pi}(s,a) = Q^{\pi}(s,a) - V^{\pi}(s)$ measures how much
better action $a$ is than the policy's average behaviour in $s$. For a
parametric policy $\pi_\theta$, the policy-gradient theorem gives an unbiased
ascent direction on $J$:

\begin{equation}
    \nabla_\theta J(\pi_\theta) \;=\;
        \mathbb{E}_{\pi_\theta}\!\left[
            \nabla_\theta \log \pi_\theta(a_t \mid s_t)\; A^{\pi_\theta}(s_t, a_t)
        \right].
    \label{eq:pg}
\end{equation}

Equation~\eqref{eq:pg} is the basis of the actor--critic methods used in this
thesis: a \emph{critic} estimates the advantage and an \emph{actor}
$\pi_\theta$ is nudged to increase the log-probability of advantageous actions.

% ---------------------------------------------------------------------------
\section{Partial Observability and Belief States}
\label{sec:bg-pomdp}

In most realistic environments the agent does not see the full state. This is
captured by the \emph{Partially Observable MDP} (POMDP)
$(\mathcal{S}, \mathcal{A}, P, R, \Omega, O, \gamma)$, which augments the MDP
with an observation space $\Omega$ and an observation model $O(o \mid s)$
\citep{kaelbling1998planning}. The agent receives only $o_t \sim O(\cdot \mid
s_t)$, which need not identify $s_t$. A policy that conditions on the current
observation alone is therefore insufficient; optimal behaviour requires
integrating the \emph{history} $\tau_t = (o_1, a_1, \dots, a_{t-1}, o_t)$.

The central object is the \emph{belief state}, the posterior over hidden states
given the history,

\begin{equation}
    b_t(s) \;=\; \Pr\!\left(s_t = s \,\middle|\, \tau_t\right),
    \label{eq:belief}
\end{equation}

updated recursively by Bayesian filtering,

\begin{equation}
    b_{t}(s') \;\propto\;
        O(o_{t} \mid s')
        \sum_{s \in \mathcal{S}} P(s' \mid s, a_{t-1})\, b_{t-1}(s).
    \label{eq:belief-update}
\end{equation}

The belief $b_t$ is a \emph{sufficient statistic} of the history for predicting
the future: it is the minimal quantity from which the reward-relevant
distribution over states can be recovered \citep{astrom1965optimal}. Crucially,
a POMDP over $\mathcal{S}$ is equivalent to a fully observable
\emph{belief MDP} over the continuous belief simplex $\Delta(\mathcal{S})$, so an
optimal Markov policy $\pi^{\star}(a \mid b)$ exists as a function of the belief
alone. This equivalence is the conceptual backbone of the thesis: \emph{the
belief is the internal variable that must be inferred, maintained, and acted
upon}, and everything we later probe or steer is an approximation to it.

Exact belief tracking is intractable for large $\mathcal{S}$. In deep RL the
belief is approximated by a learned recurrent summary of the history. A
recurrent network with hidden state $h_t \in \mathbb{R}^{d}$ implements

\begin{equation}
    h_t \;=\; f_\theta\!\left(h_{t-1},\, o_t,\, a_{t-1}\right),
    \qquad
    \pi_\theta(a_t \mid \tau_t) \;=\; \pi_\theta(a_t \mid h_t),
    \label{eq:recurrent}
\end{equation}

with $h_t$ trained end-to-end to carry exactly the history information the policy
and value heads need. We use the Gated Recurrent Unit (GRU) for $f_\theta$
\citep{cho2014learning}, whose gating stabilises credit assignment over long
horizons. Under this view $h_t$ is a \emph{learned, distributed belief}: it is
not guaranteed to equal the Bayesian posterior~\eqref{eq:belief}, but a
well-trained agent must encode within $h_t$ whatever latent structure its optimal
policy depends on. Recovering that structure from $h_t$ is the subject of
\S\ref{sec:bg-representations}.

% ---------------------------------------------------------------------------
\section{Proximal Policy Optimization}
\label{sec:bg-ppo}

We train recurrent policies with Proximal Policy Optimization (PPO)
\citep{schulman2017proximal}, a first-order policy-gradient method that is stable
and sample-efficient enough for POMDP navigation. PPO improves on the raw
gradient~\eqref{eq:pg} in two ways.

First, advantages are estimated by \emph{Generalized Advantage Estimation}
(GAE), an exponentially-weighted average of $n$-step temporal-difference errors
that trades bias for variance via $\lambda \in [0,1]$
\citep{schulman2016high}:

\begin{equation}
    \hat{A}_t \;=\; \sum_{l=0}^{\infty} (\gamma\lambda)^{l}\, \delta_{t+l},
    \qquad
    \delta_t \;=\; r_t + \gamma\, V_\phi(s_{t+1}) - V_\phi(s_t),
    \label{eq:gae}
\end{equation}

where $V_\phi$ is the learned critic. Second, PPO constrains each update so the
new policy stays close to the data-collecting policy $\pi_{\theta_{\text{old}}}$,
which prevents destructively large steps. Writing the importance ratio
$\rho_t(\theta) = \pi_\theta(a_t \mid s_t) / \pi_{\theta_{\text{old}}}(a_t \mid
s_t)$, PPO maximises the clipped surrogate

\begin{equation}
    L^{\text{CLIP}}(\theta) \;=\;
        \mathbb{E}_t\!\left[
            \min\!\Big(
                \rho_t(\theta)\, \hat{A}_t,\;
                \operatorname{clip}\!\big(\rho_t(\theta),\, 1-\epsilon,\, 1+\epsilon\big)\, \hat{A}_t
            \Big)
        \right].
    \label{eq:ppo-clip}
\end{equation}

The clip range $\epsilon$ caps the incentive to move $\rho_t$ far from $1$. The
full objective adds a value-regression term and an entropy bonus,

\begin{equation}
    L(\theta,\phi) \;=\;
        L^{\text{CLIP}}(\theta)
        \;-\; c_v\, \mathbb{E}_t\!\big[(V_\phi(s_t) - \hat{G}_t)^2\big]
        \;+\; c_e\, \mathbb{E}_t\!\big[\mathcal{H}[\pi_\theta(\cdot \mid s_t)]\big],
    \label{eq:ppo-full}
\end{equation}

where $\hat{G}_t = \hat{A}_t + V_\phi(s_t)$ is the value target and
$\mathcal{H}$ denotes entropy. The entropy coefficient $c_e$ controls
exploration; as we show in the experiments, it is decisive when the task admits
a lazy, belief-free local optimum, because escaping that optimum requires
sustained exploration of belief-conditional behaviour. For recurrent policies,
equations~\eqref{eq:gae}--\eqref{eq:ppo-full} are optimised over sequences: the
hidden state~\eqref{eq:recurrent} is carried through each rollout and reset at
episode boundaries, so gradients flow through time via backpropagation through
the unrolled GRU.

% ---------------------------------------------------------------------------
\section{Reward Shaping and Potential-Based Invariance}
\label{sec:bg-shaping}

Sparse rewards make credit assignment hard: if reward arrives only at a distant
goal, the gradient~\eqref{eq:pg} is uninformative for most of an episode.
\emph{Reward shaping} adds an auxiliary signal $F(s,a,s')$ to densify learning.
The danger is that a poorly chosen $F$ changes what is optimal, so the agent
learns to exploit the shaping instead of solving the task. \citet{ng1999policy}
show that this is avoided \emph{if and only if} the shaping is
\emph{potential-based}: for some potential function $\Phi : \mathcal{S} \to
\mathbb{R}$,

\begin{equation}
    F(s,a,s') \;=\; \gamma\, \Phi(s') - \Phi(s).
    \label{eq:pbrs}
\end{equation}

Under \eqref{eq:pbrs} the set of optimal policies is provably unchanged, because
the added term telescopes: over any trajectory the shaping contributes only
$\gamma^{T}\Phi(s_T) - \Phi(s_0)$, a quantity that depends on the endpoints, not
on the path taken. Potential-based reward shaping (PBRS) therefore accelerates
learning without biasing the solution.

In this thesis the potential is derived from a \emph{cost-to-go} oracle. Let
$\mathrm{ctg}(s)$ be the minimum number of primitive actions needed to reach the
goal from $s$, computed by a shortest-path search on the (known) map; we set

\begin{equation}
    \Phi(s) \;=\; -\,\mathrm{ctg}(s).
    \label{eq:phi-ctg}
\end{equation}

Moving closer to the goal ($\mathrm{ctg}$ decreasing) yields positive shaping,
providing a dense navigational gradient. Two properties of PBRS matter for our
analysis. First, because the shaping is invariant, it \emph{cannot} manufacture
a preference the terminal reward does not already encode; it can only make an
existing preference easier to discover. Second, its contribution to the
\emph{return difference} between two candidate endpoints is bounded by
$\lvert \Phi(s_A) - \Phi(s_B)\rvert$, which for a small
$\mathrm{ctg}$ gap can be negligible relative to a large terminal bonus. We
return to both points when explaining why shaping toward a target does not, on
its own, force a partially observable agent to condition its behaviour on the
correct latent belief.

% ---------------------------------------------------------------------------
\section{World Models and the Recurrent State-Space Model}
\label{sec:bg-worldmodel}

An alternative to learning $h_t$ purely from the RL loss is to learn an explicit
generative \emph{world model} that predicts observations and rewards, and to
train the policy inside the model's imagined rollouts. DreamerV3
\citep{hafner2023mastering,hafner2020dream} is the state-of-the-art instance and
underlies the model-based agents in this work. Its core is the \emph{Recurrent
State-Space Model} (RSSM), which factorises the latent state into a
deterministic recurrent path $h_t$ and a stochastic latent $z_t$:

\begin{align}
    h_t         &= f_\theta(h_{t-1}, z_{t-1}, a_{t-1}) && \text{(recurrent model)} \label{eq:rssm-h}\\
    \hat{z}_t   &\sim p_\theta(\hat{z}_t \mid h_t)     && \text{(prior / transition)} \label{eq:rssm-prior}\\
    z_t         &\sim q_\theta(z_t \mid h_t, x_t)      && \text{(posterior / filter)} \label{eq:rssm-post}
\end{align}

where $x_t$ is the encoded observation. The \emph{prior}
$p_\theta$~\eqref{eq:rssm-prior} predicts the next latent from dynamics alone
(enabling imagination), while the \emph{posterior}
$q_\theta$~\eqref{eq:rssm-post} corrects it using the actual observation --- an
amortised, learned version of the Bayesian filter~\eqref{eq:belief-update}. The
joint latent $s_t = (h_t, z_t)$ is the model's belief state. Decoder heads
reconstruct the observation $\hat{x}_t$ and predict reward from $s_t$, and the
whole model is trained by maximising an evidence lower bound of the form

\begin{equation}
    \mathcal{L}(\theta) \;=\;
        \underbrace{\mathbb{E}_{q}\big[\log p_\theta(x_t \mid s_t)\big]}_{\text{reconstruction}}
        \;-\; \beta\,
        \underbrace{\mathrm{KL}\!\big[\,q_\theta(z_t \mid h_t, x_t)\,\big\|\,p_\theta(\hat z_t \mid h_t)\,\big]}_{\text{dynamics / representation}}.
    \label{eq:elbo}
\end{equation}

The actor and critic are trained on trajectories \emph{imagined} by rolling the
prior forward without observations, then decoded back to interpretable
predictions. Two features of this architecture are central to the thesis.
\textbf{(i)} The belief $(h_t, z_t)$ is a single shared object consumed by
multiple heads --- the actor, the reward predictor, and the observation decoder
--- so there is no architecturally separate ``skill'' subspace. \textbf{(ii)}
Because the decoder maps the latent back to observation space, the model exposes
an explicit \emph{belief readout}: one can decode what the agent ``believes'' the
world looks like directly from $s_t$, and watch how that decoded belief changes
under intervention. This makes the world model an unusually transparent
substrate for the interpretability questions we now formalise.

% ---------------------------------------------------------------------------
\section{Representations, Probing, and the Linear Hypothesis}
\label{sec:bg-representations}

Neural agents encode task-relevant variables in \emph{distributed}
representations: a concept is not a single unit but a pattern across many units
\citep{rumelhart1986learning}. The empirically influential \emph{linear
representation hypothesis} (LRH) sharpens this: many high-level concepts are
encoded as \emph{directions} in activation space, so that the presence or value
of a concept is (approximately) an affine function of the hidden state
\citep{mikolov2013linguistic,park2024linear}. Under the LRH, a hidden state
$h \in \mathbb{R}^{d}$ decomposes as a sum of concept directions plus residual,
and the value of concept $i$ is read off by projecting onto a direction $v_i$.

The standard tool for testing this is a \emph{linear probe}: a simple classifier
trained to predict a labelled latent variable $y$ from a frozen activation $h$
\citep{alain2017understanding}. For a categorical belief label with classes
$c$, the probe is a linear--softmax map

\begin{equation}
    \hat{p}(y = c \mid h) \;=\; \operatorname{softmax}_c\!\big(W h + b\big),
    \label{eq:probe}
\end{equation}

and high held-out accuracy is taken as evidence that the information $y$ is
linearly \emph{decodable} from $h$. A key methodological caution, which this
thesis makes concrete, is that decodability is a statement about what a
\emph{downstream reader} could in principle extract --- it does \emph{not} imply
that the agent's own policy head uses that direction. In our agents we attach
such a probe as an auxiliary head trained jointly with the policy (a
``belief head''); it provides a free, continuously-monitored readout of the
latent belief without feeding back into the actor, so that decodability and use
can be measured separately. When a concept is linearly represented, a canonical
estimate of its direction is the \emph{difference of class means} in activation
space,

\begin{equation}
    v \;=\; \mu_{+} - \mu_{-}, \qquad
    \mu_{\pm} \;=\; \frac{1}{|\mathcal{D}_{\pm}|} \sum_{h \in \mathcal{D}_{\pm}} h,
    \label{eq:diffmeans}
\end{equation}

where $\mathcal{D}_{+}$ and $\mathcal{D}_{-}$ are activations collected under the
two values of the concept \citep{marks2024geometry}. The vector $v$ points from
one class centroid to the other and is the workhorse direction for the
interventions in the next section. The LRH is, however, only approximate:
\emph{superposition} lets networks pack more features than dimensions by
overlapping them in non-orthogonal directions \citep{elhage2022superposition},
and some features occupy multi-dimensional subspaces rather than single lines.
This non-orthogonality is precisely what makes targeted intervention leaky, a
theme we develop in \S\ref{sec:bg-belief}.

% ---------------------------------------------------------------------------
\section{Activation Steering}
\label{sec:bg-steering}

\emph{Activation steering} intervenes on a model's computation at inference time
by editing its hidden states along a chosen direction, rather than retraining
weights. Given a steering direction $v$ (often estimated by
\eqref{eq:diffmeans}), a layer activation $h$, and a scale $\alpha$, the
canonical edit is additive:

\begin{equation}
    h' \;=\; h + \alpha\, v.
    \label{eq:steer-add}
\end{equation}

This is \emph{activation addition} (ActAdd), which injects a contrast vector at
inference to push behaviour toward or away from a target concept
\citep{turner2023activation}. \emph{Contrastive Activation Addition} (CAA)
constructs $v$ as the mean difference of activations over matched positive and
negative prompt pairs \citep{rimsky2024steering}, i.e.\ a structured instance of
\eqref{eq:diffmeans}. A complementary operation is directional \emph{ablation},
which removes a concept by projecting it out of the activation,

\begin{equation}
    h' \;=\; h - (\hat{v}^{\top} h)\, \hat{v}, \qquad \hat{v} = v / \lVert v \rVert,
    \label{eq:steer-ablate}
\end{equation}

as used to suppress the refusal direction in language models
\citep{arditi2024refusal}. Finer-grained variants steer individual features
discovered by \emph{sparse autoencoders}, which decompose $h$ into a sparse,
ideally monosemantic dictionary and amplify or suppress single entries
\citep{bricken2023monosemanticity,cunningham2024sparse}. In this thesis we apply
\eqref{eq:steer-add}--\eqref{eq:steer-ablate} to the recurrent belief carrier
$h_t$ of an RL agent: the direction $v$ is a difference-of-means over the latent
belief (e.g.\ activations on maps of one category minus another), and it is
injected over a window of timesteps while the agent acts. The behavioural
consequence --- whether the intervention changes the agent's downstream choice
--- is then read out from its actions, and the belief consequence from the
decoded world model or belief probe.

Every method in this family shares an implicit assumption: that editing along
$v$ leaves the computations served by other directions untouched. Because
representations are entangled and non-orthogonal
(\S\ref{sec:bg-representations}), any edit with a nonzero projection onto other
functional directions perturbs them too. Whether a given belief direction is
in fact \emph{used} by the policy, and whether steering it produces a clean
behavioural change or leaks into unrelated capabilities, are empirical questions
--- and the ones this thesis is built to answer.

% ---------------------------------------------------------------------------
\section{Belief, Skill, and the Decoupling Question}
\label{sec:bg-belief}

We can now state the conceptual frame precisely. We distinguish two readouts of a
shared latent state $z$ (either the GRU state $h_t$ or the RSSM latent
$(h_t,z_t)$). The \emph{belief readout} $B(z)$ is the model's inferred
distribution over hidden world structure --- e.g.\ the decoded observation of a
world model, or the belief probe~\eqref{eq:probe}. The \emph{policy readout}
$\pi(z)$ is the distribution over actions. A latent variable is
\emph{behaviourally load-bearing} when the policy's optimal action depends on it;
it is \emph{merely decodable} when $B(z)$ recovers it but $\pi(z)$ does not use
it.

Localised steering implicitly assumes these are separable: that one can edit the
belief without disturbing the policy, or edit the policy's target behaviour
without disturbing the belief. Following the diagnostic used in this work, an
intervention $\delta$ applied to $z$ is \emph{belief-preserving} iff it moves the
policy while leaving the belief readout intact:

\begin{equation}
    \big\lVert B(z + \delta) - B(z) \big\rVert < \varepsilon
    \quad\text{and}\quad
    \big\lVert \pi(z + \delta) - \pi(z) \big\rVert > \eta,
    \label{eq:belief-invariance}
\end{equation}

for pre-specified tolerances $\varepsilon$ (permitted belief drift) and $\eta$
(required behavioural change). Symmetrically, a behaviour-targeting edit is
\emph{clean} only if it does not silently shift $B(z)$. When belief and policy
readouts load onto overlapping directions of a shared latent, no strictly linear
edit can satisfy \eqref{eq:belief-invariance} exactly; the severity of the
resulting cross-talk is a property of the representational geometry, not of the
method. This is why architectures with an explicit belief decoder are valuable:
they let $B(z)$ be measured directly, turning \eqref{eq:belief-invariance} from a
philosophical distinction into a checkable quantity.

The environments and agents in the following chapters are designed to make this
question concrete and controllable. The agent must infer a latent map property
(the belief) from a partial history and select among behaviours accordingly (the
skill), while an auxiliary belief head~\eqref{eq:probe} and, for the world-model
agents, an observation decoder provide independent readouts of $B(z)$. This
setup lets us ask, and answer empirically, the two questions that motivate the
thesis: \emph{(i)} when is a decodable belief actually used by the policy, and
\emph{(ii)} does steering the belief direction change behaviour as intended, or
does it drift the belief while leaving behaviour untouched --- or vice versa.
.
% ============================================================================

\chapter{Background}
\label{ch:background}

This chapter develops the theory required to follow the rest of the thesis. We
proceed from the general framework of sequential decision making under
uncertainty (\S\ref{sec:bg-rl}--\S\ref{sec:bg-pomdp}), through the specific
learning algorithms used to train our agents (\S\ref{sec:bg-ppo}) and the reward
machinery that shapes them (\S\ref{sec:bg-shaping}), to the model-based world
models that motivate the belief-centric view (\S\ref{sec:bg-worldmodel}).
We then introduce the interpretability tools that this work rests on: the
representational hypotheses that make internal states legible
(\S\ref{sec:bg-representations}), the family of activation-steering interventions
we study (\S\ref{sec:bg-steering}), and finally the belief/behaviour distinction
that frames our central question (\S\ref{sec:bg-belief}).

% ---------------------------------------------------------------------------
\section{Reinforcement Learning and Markov Decision Processes}
\label{sec:bg-rl}

Reinforcement learning (RL) formalises an agent that learns to act by trial and
error so as to maximise a scalar reward \citep{sutton2018reinforcement}. The
standard model is the \emph{Markov Decision Process} (MDP), a tuple
$\mathcal{M} = (\mathcal{S}, \mathcal{A}, P, R, \gamma)$ with state space
$\mathcal{S}$, action space $\mathcal{A}$, transition kernel
$P(s' \mid s, a)$, reward function $R(s,a)$, and discount factor
$\gamma \in [0,1)$. At each step $t$ the agent observes state $s_t$, samples an
action $a_t \sim \pi(\cdot \mid s_t)$ from its \emph{policy} $\pi$, receives
reward $r_t = R(s_t, a_t)$, and transitions to $s_{t+1} \sim P(\cdot \mid s_t,
a_t)$. The objective is to maximise the expected discounted return

\begin{equation}
    G_t \;=\; \sum_{k=0}^{\infty} \gamma^{k}\, r_{t+k},
    \qquad
    J(\pi) \;=\; \mathbb{E}_{\pi}\!\left[G_0\right].
    \label{eq:return}
\end{equation}

The discount $\gamma$ trades off immediate against future reward and keeps the
sum finite. Two value functions summarise long-run desirability: the
\emph{state value} $V^{\pi}(s) = \mathbb{E}_{\pi}[G_t \mid s_t = s]$ and the
\emph{action value} $Q^{\pi}(s,a) = \mathbb{E}_{\pi}[G_t \mid s_t = s, a_t = a]$.
They satisfy the Bellman consistency equations, e.g.

\begin{equation}
    V^{\pi}(s) \;=\; \mathbb{E}_{a\sim\pi,\, s'\sim P}
        \big[\, R(s,a) + \gamma\, V^{\pi}(s') \,\big].
    \label{eq:bellman}
\end{equation}

The \emph{advantage} $A^{\pi}(s,a) = Q^{\pi}(s,a) - V^{\pi}(s)$ measures how much
better action $a$ is than the policy's average behaviour in $s$. For a
parametric policy $\pi_\theta$, the policy-gradient theorem gives an unbiased
ascent direction on $J$:

\begin{equation}
    \nabla_\theta J(\pi_\theta) \;=\;
        \mathbb{E}_{\pi_\theta}\!\left[
            \nabla_\theta \log \pi_\theta(a_t \mid s_t)\; A^{\pi_\theta}(s_t, a_t)
        \right].
    \label{eq:pg}
\end{equation}

Equation~\eqref{eq:pg} is the basis of the actor--critic methods used in this
thesis: a \emph{critic} estimates the advantage and an \emph{actor}
$\pi_\theta$ is nudged to increase the log-probability of advantageous actions.

% ---------------------------------------------------------------------------
\section{Partial Observability and Belief States}
\label{sec:bg-pomdp}

In most realistic environments the agent does not see the full state. This is
captured by the \emph{Partially Observable MDP} (POMDP)
$(\mathcal{S}, \mathcal{A}, P, R, \Omega, O, \gamma)$, which augments the MDP
with an observation space $\Omega$ and an observation model $O(o \mid s)$
\citep{kaelbling1998planning}. The agent receives only $o_t \sim O(\cdot \mid
s_t)$, which need not identify $s_t$. A policy that conditions on the current
observation alone is therefore insufficient; optimal behaviour requires
integrating the \emph{history} $\tau_t = (o_1, a_1, \dots, a_{t-1}, o_t)$.

The central object is the \emph{belief state}, the posterior over hidden states
given the history,

\begin{equation}
    b_t(s) \;=\; \Pr\!\left(s_t = s \,\middle|\, \tau_t\right),
    \label{eq:belief}
\end{equation}

updated recursively by Bayesian filtering,

\begin{equation}
    b_{t}(s') \;\propto\;
        O(o_{t} \mid s')
        \sum_{s \in \mathcal{S}} P(s' \mid s, a_{t-1})\, b_{t-1}(s).
    \label{eq:belief-update}
\end{equation}

The belief $b_t$ is a \emph{sufficient statistic} of the history for predicting
the future: it is the minimal quantity from which the reward-relevant
distribution over states can be recovered \citep{astrom1965optimal}. Crucially,
a POMDP over $\mathcal{S}$ is equivalent to a fully observable
\emph{belief MDP} over the continuous belief simplex $\Delta(\mathcal{S})$, so an
optimal Markov policy $\pi^{\star}(a \mid b)$ exists as a function of the belief
alone. This equivalence is the conceptual backbone of the thesis: \emph{the
belief is the internal variable that must be inferred, maintained, and acted
upon}, and everything we later probe or steer is an approximation to it.

Exact belief tracking is intractable for large $\mathcal{S}$. In deep RL the
belief is approximated by a learned recurrent summary of the history. A
recurrent network with hidden state $h_t \in \mathbb{R}^{d}$ implements

\begin{equation}
    h_t \;=\; f_\theta\!\left(h_{t-1},\, o_t,\, a_{t-1}\right),
    \qquad
    \pi_\theta(a_t \mid \tau_t) \;=\; \pi_\theta(a_t \mid h_t),
    \label{eq:recurrent}
\end{equation}

with $h_t$ trained end-to-end to carry exactly the history information the policy
and value heads need. We use the Gated Recurrent Unit (GRU) for $f_\theta$
\citep{cho2014learning}, whose gating stabilises credit assignment over long
horizons. Under this view $h_t$ is a \emph{learned, distributed belief}: it is
not guaranteed to equal the Bayesian posterior~\eqref{eq:belief}, but a
well-trained agent must encode within $h_t$ whatever latent structure its optimal
policy depends on. Recovering that structure from $h_t$ is the subject of
\S\ref{sec:bg-representations}.

% ---------------------------------------------------------------------------
\section{Proximal Policy Optimization}
\label{sec:bg-ppo}

We train recurrent policies with Proximal Policy Optimization (PPO)
\citep{schulman2017proximal}, a first-order policy-gradient method that is stable
and sample-efficient enough for POMDP navigation. PPO improves on the raw
gradient~\eqref{eq:pg} in two ways.

First, advantages are estimated by \emph{Generalized Advantage Estimation}
(GAE), an exponentially-weighted average of $n$-step temporal-difference errors
that trades bias for variance via $\lambda \in [0,1]$
\citep{schulman2016high}:

\begin{equation}
    \hat{A}_t \;=\; \sum_{l=0}^{\infty} (\gamma\lambda)^{l}\, \delta_{t+l},
    \qquad
    \delta_t \;=\; r_t + \gamma\, V_\phi(s_{t+1}) - V_\phi(s_t),
    \label{eq:gae}
\end{equation}

where $V_\phi$ is the learned critic. Second, PPO constrains each update so the
new policy stays close to the data-collecting policy $\pi_{\theta_{\text{old}}}$,
which prevents destructively large steps. Writing the importance ratio
$\rho_t(\theta) = \pi_\theta(a_t \mid s_t) / \pi_{\theta_{\text{old}}}(a_t \mid
s_t)$, PPO maximises the clipped surrogate

\begin{equation}
    L^{\text{CLIP}}(\theta) \;=\;
        \mathbb{E}_t\!\left[
            \min\!\Big(
                \rho_t(\theta)\, \hat{A}_t,\;
                \operatorname{clip}\!\big(\rho_t(\theta),\, 1-\epsilon,\, 1+\epsilon\big)\, \hat{A}_t
            \Big)
        \right].
    \label{eq:ppo-clip}
\end{equation}

The clip range $\epsilon$ caps the incentive to move $\rho_t$ far from $1$. The
full objective adds a value-regression term and an entropy bonus,

\begin{equation}
    L(\theta,\phi) \;=\;
        L^{\text{CLIP}}(\theta)
        \;-\; c_v\, \mathbb{E}_t\!\big[(V_\phi(s_t) - \hat{G}_t)^2\big]
        \;+\; c_e\, \mathbb{E}_t\!\big[\mathcal{H}[\pi_\theta(\cdot \mid s_t)]\big],
    \label{eq:ppo-full}
\end{equation}

where $\hat{G}_t = \hat{A}_t + V_\phi(s_t)$ is the value target and
$\mathcal{H}$ denotes entropy. The entropy coefficient $c_e$ controls
exploration; as we show in the experiments, it is decisive when the task admits
a lazy, belief-free local optimum, because escaping that optimum requires
sustained exploration of belief-conditional behaviour. For recurrent policies,
equations~\eqref{eq:gae}--\eqref{eq:ppo-full} are optimised over sequences: the
hidden state~\eqref{eq:recurrent} is carried through each rollout and reset at
episode boundaries, so gradients flow through time via backpropagation through
the unrolled GRU.

% ---------------------------------------------------------------------------
\section{Reward Shaping and Potential-Based Invariance}
\label{sec:bg-shaping}

Sparse rewards make credit assignment hard: if reward arrives only at a distant
goal, the gradient~\eqref{eq:pg} is uninformative for most of an episode.
\emph{Reward shaping} adds an auxiliary signal $F(s,a,s')$ to densify learning.
The danger is that a poorly chosen $F$ changes what is optimal, so the agent
learns to exploit the shaping instead of solving the task. \citet{ng1999policy}
show that this is avoided \emph{if and only if} the shaping is
\emph{potential-based}: for some potential function $\Phi : \mathcal{S} \to
\mathbb{R}$,

\begin{equation}
    F(s,a,s') \;=\; \gamma\, \Phi(s') - \Phi(s).
    \label{eq:pbrs}
\end{equation}

Under \eqref{eq:pbrs} the set of optimal policies is provably unchanged, because
the added term telescopes: over any trajectory the shaping contributes only
$\gamma^{T}\Phi(s_T) - \Phi(s_0)$, a quantity that depends on the endpoints, not
on the path taken. Potential-based reward shaping (PBRS) therefore accelerates
learning without biasing the solution.

In this thesis the potential is derived from a \emph{cost-to-go} oracle. Let
$\mathrm{ctg}(s)$ be the minimum number of primitive actions needed to reach the
goal from $s$, computed by a shortest-path search on the (known) map; we set

\begin{equation}
    \Phi(s) \;=\; -\,\mathrm{ctg}(s).
    \label{eq:phi-ctg}
\end{equation}

Moving closer to the goal ($\mathrm{ctg}$ decreasing) yields positive shaping,
providing a dense navigational gradient. Two properties of PBRS matter for our
analysis. First, because the shaping is invariant, it \emph{cannot} manufacture
a preference the terminal reward does not already encode; it can only make an
existing preference easier to discover. Second, its contribution to the
\emph{return difference} between two candidate endpoints is bounded by
$\lvert \Phi(s_A) - \Phi(s_B)\rvert$, which for a small
$\mathrm{ctg}$ gap can be negligible relative to a large terminal bonus. We
return to both points when explaining why shaping toward a target does not, on
its own, force a partially observable agent to condition its behaviour on the
correct latent belief.

% ---------------------------------------------------------------------------
\section{World Models and the Recurrent State-Space Model}
\label{sec:bg-worldmodel}

An alternative to learning $h_t$ purely from the RL loss is to learn an explicit
generative \emph{world model} that predicts observations and rewards, and to
train the policy inside the model's imagined rollouts. DreamerV3
\citep{hafner2023mastering,hafner2020dream} is the state-of-the-art instance and
underlies the model-based agents in this work. Its core is the \emph{Recurrent
State-Space Model} (RSSM), which factorises the latent state into a
deterministic recurrent path $h_t$ and a stochastic latent $z_t$:

\begin{align}
    h_t         &= f_\theta(h_{t-1}, z_{t-1}, a_{t-1}) && \text{(recurrent model)} \label{eq:rssm-h}\\
    \hat{z}_t   &\sim p_\theta(\hat{z}_t \mid h_t)     && \text{(prior / transition)} \label{eq:rssm-prior}\\
    z_t         &\sim q_\theta(z_t \mid h_t, x_t)      && \text{(posterior / filter)} \label{eq:rssm-post}
\end{align}

where $x_t$ is the encoded observation. The \emph{prior}
$p_\theta$~\eqref{eq:rssm-prior} predicts the next latent from dynamics alone
(enabling imagination), while the \emph{posterior}
$q_\theta$~\eqref{eq:rssm-post} corrects it using the actual observation --- an
amortised, learned version of the Bayesian filter~\eqref{eq:belief-update}. The
joint latent $s_t = (h_t, z_t)$ is the model's belief state. Decoder heads
reconstruct the observation $\hat{x}_t$ and predict reward from $s_t$, and the
whole model is trained by maximising an evidence lower bound of the form

\begin{equation}
    \mathcal{L}(\theta) \;=\;
        \underbrace{\mathbb{E}_{q}\big[\log p_\theta(x_t \mid s_t)\big]}_{\text{reconstruction}}
        \;-\; \beta\,
        \underbrace{\mathrm{KL}\!\big[\,q_\theta(z_t \mid h_t, x_t)\,\big\|\,p_\theta(\hat z_t \mid h_t)\,\big]}_{\text{dynamics / representation}}.
    \label{eq:elbo}
\end{equation}

The actor and critic are trained on trajectories \emph{imagined} by rolling the
prior forward without observations, then decoded back to interpretable
predictions. Two features of this architecture are central to the thesis.
\textbf{(i)} The belief $(h_t, z_t)$ is a single shared object consumed by
multiple heads --- the actor, the reward predictor, and the observation decoder
--- so there is no architecturally separate ``skill'' subspace. \textbf{(ii)}
Because the decoder maps the latent back to observation space, the model exposes
an explicit \emph{belief readout}: one can decode what the agent ``believes'' the
world looks like directly from $s_t$, and watch how that decoded belief changes
under intervention. This makes the world model an unusually transparent
substrate for the interpretability questions we now formalise.

% ---------------------------------------------------------------------------
\section{Representations, Probing, and the Linear Hypothesis}
\label{sec:bg-representations}

Neural agents encode task-relevant variables in \emph{distributed}
representations: a concept is not a single unit but a pattern across many units
\citep{rumelhart1986learning}. The empirically influential \emph{linear
representation hypothesis} (LRH) sharpens this: many high-level concepts are
encoded as \emph{directions} in activation space, so that the presence or value
of a concept is (approximately) an affine function of the hidden state
\citep{mikolov2013linguistic,park2024linear}. Under the LRH, a hidden state
$h \in \mathbb{R}^{d}$ decomposes as a sum of concept directions plus residual,
and the value of concept $i$ is read off by projecting onto a direction $v_i$.

The standard tool for testing this is a \emph{linear probe}: a simple classifier
trained to predict a labelled latent variable $y$ from a frozen activation $h$
\citep{alain2017understanding}. For a categorical belief label with classes
$c$, the probe is a linear--softmax map

\begin{equation}
    \hat{p}(y = c \mid h) \;=\; \operatorname{softmax}_c\!\big(W h + b\big),
    \label{eq:probe}
\end{equation}

and high held-out accuracy is taken as evidence that the information $y$ is
linearly \emph{decodable} from $h$. A key methodological caution, which this
thesis makes concrete, is that decodability is a statement about what a
\emph{downstream reader} could in principle extract --- it does \emph{not} imply
that the agent's own policy head uses that direction. In our agents we attach
such a probe as an auxiliary head trained jointly with the policy (a
``belief head''); it provides a free, continuously-monitored readout of the
latent belief without feeding back into the actor, so that decodability and use
can be measured separately. When a concept is linearly represented, a canonical
estimate of its direction is the \emph{difference of class means} in activation
space,

\begin{equation}
    v \;=\; \mu_{+} - \mu_{-}, \qquad
    \mu_{\pm} \;=\; \frac{1}{|\mathcal{D}_{\pm}|} \sum_{h \in \mathcal{D}_{\pm}} h,
    \label{eq:diffmeans}
\end{equation}

where $\mathcal{D}_{+}$ and $\mathcal{D}_{-}$ are activations collected under the
two values of the concept \citep{marks2024geometry}. The vector $v$ points from
one class centroid to the other and is the workhorse direction for the
interventions in the next section. The LRH is, however, only approximate:
\emph{superposition} lets networks pack more features than dimensions by
overlapping them in non-orthogonal directions \citep{elhage2022superposition},
and some features occupy multi-dimensional subspaces rather than single lines.
This non-orthogonality is precisely what makes targeted intervention leaky, a
theme we develop in \S\ref{sec:bg-belief}.

% ---------------------------------------------------------------------------
\section{Activation Steering}
\label{sec:bg-steering}

\emph{Activation steering} intervenes on a model's computation at inference time
by editing its hidden states along a chosen direction, rather than retraining
weights. Given a steering direction $v$ (often estimated by
\eqref{eq:diffmeans}), a layer activation $h$, and a scale $\alpha$, the
canonical edit is additive:

\begin{equation}
    h' \;=\; h + \alpha\, v.
    \label{eq:steer-add}
\end{equation}

This is \emph{activation addition} (ActAdd), which injects a contrast vector at
inference to push behaviour toward or away from a target concept
\citep{turner2023activation}. \emph{Contrastive Activation Addition} (CAA)
constructs $v$ as the mean difference of activations over matched positive and
negative prompt pairs \citep{rimsky2024steering}, i.e.\ a structured instance of
\eqref{eq:diffmeans}. A complementary operation is directional \emph{ablation},
which removes a concept by projecting it out of the activation,

\begin{equation}
    h' \;=\; h - (\hat{v}^{\top} h)\, \hat{v}, \qquad \hat{v} = v / \lVert v \rVert,
    \label{eq:steer-ablate}
\end{equation}

as used to suppress the refusal direction in language models
\citep{arditi2024refusal}. Finer-grained variants steer individual features
discovered by \emph{sparse autoencoders}, which decompose $h$ into a sparse,
ideally monosemantic dictionary and amplify or suppress single entries
\citep{bricken2023monosemanticity,cunningham2024sparse}. In this thesis we apply
\eqref{eq:steer-add}--\eqref{eq:steer-ablate} to the recurrent belief carrier
$h_t$ of an RL agent: the direction $v$ is a difference-of-means over the latent
belief (e.g.\ activations on maps of one category minus another), and it is
injected over a window of timesteps while the agent acts. The behavioural
consequence --- whether the intervention changes the agent's downstream choice
--- is then read out from its actions, and the belief consequence from the
decoded world model or belief probe.

Every method in this family shares an implicit assumption: that editing along
$v$ leaves the computations served by other directions untouched. Because
representations are entangled and non-orthogonal
(\S\ref{sec:bg-representations}), any edit with a nonzero projection onto other
functional directions perturbs them too. Whether a given belief direction is
in fact \emph{used} by the policy, and whether steering it produces a clean
behavioural change or leaks into unrelated capabilities, are empirical questions
--- and the ones this thesis is built to answer.

% ---------------------------------------------------------------------------
\section{Belief, Skill, and the Decoupling Question}
\label{sec:bg-belief}

We can now state the conceptual frame precisely. We distinguish two readouts of a
shared latent state $z$ (either the GRU state $h_t$ or the RSSM latent
$(h_t,z_t)$). The \emph{belief readout} $B(z)$ is the model's inferred
distribution over hidden world structure --- e.g.\ the decoded observation of a
world model, or the belief probe~\eqref{eq:probe}. The \emph{policy readout}
$\pi(z)$ is the distribution over actions. A latent variable is
\emph{behaviourally load-bearing} when the policy's optimal action depends on it;
it is \emph{merely decodable} when $B(z)$ recovers it but $\pi(z)$ does not use
it.

Localised steering implicitly assumes these are separable: that one can edit the
belief without disturbing the policy, or edit the policy's target behaviour
without disturbing the belief. Following the diagnostic used in this work, an
intervention $\delta$ applied to $z$ is \emph{belief-preserving} iff it moves the
policy while leaving the belief readout intact:

\begin{equation}
    \big\lVert B(z + \delta) - B(z) \big\rVert < \varepsilon
    \quad\text{and}\quad
    \big\lVert \pi(z + \delta) - \pi(z) \big\rVert > \eta,
    \label{eq:belief-invariance}
\end{equation}

for pre-specified tolerances $\varepsilon$ (permitted belief drift) and $\eta$
(required behavioural change). Symmetrically, a behaviour-targeting edit is
\emph{clean} only if it does not silently shift $B(z)$. When belief and policy
readouts load onto overlapping directions of a shared latent, no strictly linear
edit can satisfy \eqref{eq:belief-invariance} exactly; the severity of the
resulting cross-talk is a property of the representational geometry, not of the
method. This is why architectures with an explicit belief decoder are valuable:
they let $B(z)$ be measured directly, turning \eqref{eq:belief-invariance} from a
philosophical distinction into a checkable quantity.

The environments and agents in the following chapters are designed to make this
question concrete and controllable. The agent must infer a latent map property
(the belief) from a partial history and select among behaviours accordingly (the
skill), while an auxiliary belief head~\eqref{eq:probe} and, for the world-model
agents, an observation decoder provide independent readouts of $B(z)$. This
setup lets us ask, and answer empirically, the two questions that motivate the
thesis: \emph{(i)} when is a decodable belief actually used by the policy, and
\emph{(ii)} does steering the belief direction change behaviour as intended, or
does it drift the belief while leaving behaviour untouched --- or vice versa.
